\documentclass[journal]{IEEEtran}

\usepackage{amsmath,amssymb,bm}
\usepackage{cite}
\usepackage{booktabs}
\usepackage{array}
\usepackage{url}
\usepackage{enumitem}

\interdisplaylinepenalty=2500

\newcommand{\gsid}{\mathrm{GeoSemID}}
\newcommand{\Experts}{\mathcal{M}}
\newcommand{\POIs}{\mathcal{P}}
\newcommand{\Users}{\mathcal{U}}
\newcommand{\Train}{\mathcal{D}_{\mathrm{tr}}}
\newcommand{\HitTrain}{\mathcal{D}_{\mathrm{hit}}}
\newcommand{\Ind}{\mathbb{I}}

\title{Constrained Geo-Semantic Identifiers and Utility-Supervised Dynamic Multi-Expert Retrieval for Next-POI Prediction}

\author{Anonymous~Authors%
\thanks{Working draft prepared for IEEE Transactions on Intelligent Transportation Systems.}
}

\begin{document}

\maketitle

\begin{abstract}
Next point-of-interest (POI) prediction requires identifying a user's next
destination from a large and heterogeneous location catalog. Existing methods
often represent POIs as isolated identifiers and rely on a single retrieval
signal, which limits candidate coverage for sparse trajectories and long-tail
locations. We propose a GeoSemID framework that combines constrained
geo-semantic identifiers with utility-supervised dynamic multi-expert
retrieval. GeoSemID augments each POI with deterministic coarse- and
fine-grained geographic codes, a normalized category code, and a validated
semantic code generated offline from POI metadata. Spatial, transition,
temporal, and semantic experts generate complementary candidates; count-based
experts use reliability-aware hierarchical backoff, while spatial and semantic
experts use explicitly calibrated evidence fusion. A trajectory-conditioned
router is supervised only by training examples whose target is retrieved by at
least one expert, and it fuses rank-percentile evidence that is retained by a
candidate-constrained language-model ranker. The ranker explicitly scores only
valid candidate IDs, thereby separating candidate coverage from conditional
ranking and avoiding target leakage. Preliminary evaluations on the CA and NYK
datasets support the effectiveness of the GeoSemID representation and dynamic
retrieval mechanism; complete benchmark and ablation results will be reported
after all experiments are finalized.
\end{abstract}

\begin{IEEEkeywords}
Next point-of-interest prediction, trajectory modeling, semantic identifiers,
multi-expert retrieval, dynamic routing, large language models.
\end{IEEEkeywords}

\section{Problem Formulation}
\label{sec:problem}

\subsection{Next-POI Prediction}

Let $\Users=\{u_1,\ldots,u_{|\Users|}\}$ and
$\POIs=\{p_1,\ldots,p_{|\POIs|}\}$ denote the user and POI sets.
Each POI $p$ is associated with metadata
\begin{equation}
\bm{m}_{p}
=
\langle n_p,c_p,\bm{g}_p,\bm{a}_p\rangle,
\label{eq:poi_metadata}
\end{equation}
where $n_p$ is the POI name, $c_p$ is its functional category,
$\bm{g}_p$ contains coordinates and available administrative information,
and $\bm{a}_p$ denotes other attributes.

A check-in of user $u$ at sequence position $i$ is
\begin{equation}
x_{u,i}=\langle p_{u,i},t_{u,i}\rangle.
\label{eq:checkin}
\end{equation}
For a sample whose target is $p_u^{*}=p_{u,n_u}$, the observable trajectory is
$\{x_{u,1},\ldots,x_{u,n_u-1}\}$.
Given a recent-window length $L$, it is divided into
\begin{align}
\mathcal{T}_{u}^{H}
&=
\{x_{u,1},\ldots,x_{u,n_u-L-1}\},
\label{eq:historical_trajectory}\\
\mathcal{T}_{u}^{C}
&=
\{x_{u,n_u-L},\ldots,x_{u,n_u-1}\},
\label{eq:current_trajectory}
\end{align}
where $\mathcal{T}_{u}^{H}$ characterizes long-term mobility preference and
$\mathcal{T}_{u}^{C}$ describes the immediate movement context.

Because the complete POI catalog contains many irrelevant alternatives, we
formulate next-POI prediction as two coupled tasks.
First, a retrieval model constructs a compact candidate set
$\mathcal{C}_u\subset\POIs$ and should maximize
\begin{equation}
\Pr(p_u^{*}\in\mathcal{C}_u
\mid
\mathcal{T}_{u}^{H},\mathcal{T}_{u}^{C},\{\bm{m}_p\}_{p\in\POIs}).
\label{eq:coverage_objective}
\end{equation}
Second, a ranker predicts
\begin{equation}
\widehat{p}_{u}
=
\arg\max_{p\in\mathcal{C}_u}
P_{\psi}
\left(
p\mid
\mathcal{T}_{u}^{H},
\mathcal{T}_{u}^{C},
\mathcal{C}_u,
\gsid(p),
\bm{\xi}_{u,p}
\right),
\label{eq:final_prediction}
\end{equation}
where $\gsid(p)$ is the structured identifier of candidate $p$ and
$\bm{\xi}_{u,p}$ records expert retrieval evidence.
If $p_u^{*}\notin\mathcal{C}_u$, the end-to-end prediction is necessarily
incorrect; therefore candidate coverage and conditional ranking quality must
be evaluated separately.

\subsection{Constrained Geo-Semantic Identification}

An atomic POI ID preserves exact identity but does not expose reusable
geographic, functional, or semantic relations.
For each POI, we define
\begin{equation}
\gsid(p)
=
[
z_{p}^{\mathrm{cg}},
z_{p}^{\mathrm{fg}},
z_{p}^{\mathrm{cat}},
z_{p}^{\mathrm{sem}}
],
\label{eq:geosemid_definition}
\end{equation}
where $z_{p}^{\mathrm{cg}}$ and $z_{p}^{\mathrm{fg}}$ denote coarse- and
fine-grained geographic cells, $z_{p}^{\mathrm{cat}}$ is a normalized
functional code, and $z_{p}^{\mathrm{sem}}$ is a constrained semantic code.

The first three components are deterministic functions of metadata:
\begin{align}
z_{p}^{\mathrm{cg}}&=\rho_{\mathrm{cg}}(\bm{g}_p),&
z_{p}^{\mathrm{fg}}&=\rho_{\mathrm{fg}}(\bm{g}_p),&
z_{p}^{\mathrm{cat}}&=\rho_{\mathrm{cat}}(c_p).
\label{eq:deterministic_codes}
\end{align}
The semantic component is produced offline by a fine-tuned LLM generator:
\begin{equation}
z_{p}^{\mathrm{sem}}
=
\operatorname{Validate}
\left[
G_{\phi}\!\left(\mathcal{I}(\bm{m}_p)\right)
\right],
\label{eq:semantic_generation}
\end{equation}
where $\mathcal{I}(\cdot)$ serializes POI metadata under a fixed prompt,
$G_{\phi}$ is the generator, and $\operatorname{Validate}(\cdot)$ enforces
the predefined grammar and vocabulary.
The generator is not described as an autonomous agent: it performs a single
constrained metadata-to-code mapping.

The semantic-code target has the fixed form
\begin{equation}
y_p^{\mathrm{sem}}
=
\langle
y_p^{\mathrm{intent}},
y_p^{\mathrm{attribute}},
y_p^{\mathrm{context}}
\rangle
\in
\mathcal{V}_{i}\times\mathcal{V}_{a}\times\mathcal{V}_{c}.
\label{eq:semantic_schema}
\end{equation}
The category field is deliberately excluded from this code: it is represented
only by $z_p^{\mathrm{cat}}$. The semantic fields instead encode visit intent,
distinctive service attributes, and contextual usage beyond the source
taxonomy. The label-construction protocol, vocabulary sizes, and quality-control
statistics are reported with the experimental settings.

\subsection{Context-Adaptive Multi-Expert Retrieval}

Let
\begin{equation}
\Experts=\{\mathrm{sp},\mathrm{tr},\mathrm{tp},\mathrm{se}\}
\label{eq:expert_set}
\end{equation}
denote spatial, transition, temporal, and semantic experts.
Each expert scores the complete catalog during offline training and returns
its top-$K_0$ list during inference:
\begin{equation}
\mathcal{C}_{u}^{(m)}
=
\operatorname{TopK}_{K_0}
\{r_{u}^{(m)}(p):p\in\POIs\}.
\label{eq:expert_candidate}
\end{equation}
A sample-level router estimates $\bm{\pi}_u=\{\pi_u^{(m)}\}_{m\in\Experts}$
from observable trajectory context only.
The router fuses rank-normalized expert evidence into the final
$\mathcal{C}_u$.
Training-target ranks are used only to construct router supervision on the
training split and are never inputs to the router at validation or test time.

\section{Methodology}
\label{sec:method}

\subsection{Framework Overview}

The framework contains three coupled stages.
First, GeoSemID converts heterogeneous POI metadata into a reusable hierarchy
while preserving the original POI ID for exact identity.
Second, transition and temporal retrievers use count-reliability-aware
hierarchical backoff, whereas spatial and semantic retrievers use calibrated
hierarchical evidence; a target-rank-supervised router dynamically selects
their relative importance for each sample.
Third, a candidate-constrained LLM scores only valid candidate identifiers and
retains expert evidence in its input.

The novelty boundary is therefore not ``using Qwen'' or merely combining four
retrievers.
The intended technical contributions are:
(i) a constrained, reusable geo-semantic code that enables structured backoff
for sparse POIs;
(ii) sample-level utility supervision for the dynamic router rather than
fixed fusion; and
(iii) evidence preservation from retrieval to candidate ranking.

\subsection{Offline GeoSemID Construction}
\label{subsec:geosemid}

\subsubsection{Semantic-Code Supervision and Data Isolation}

The semantic generator is trained on
\begin{equation}
\mathcal{D}_{\mathrm{sem}}
=
\{
(\mathcal{I}(\bm{m}_p),y_p^{\mathrm{sem}})
\},
\label{eq:semantic_dataset}
\end{equation}
using token-level negative log-likelihood
\begin{equation}
\mathcal{L}_{\mathrm{sem}}
=
-
\sum_{(\bm{x},\bm{y})\in\mathcal{D}_{\mathrm{sem}}}
\sum_{\ell=1}^{|\bm{y}|}
\log P_{\phi}(y_{\ell}\mid\bm{x},y_{<\ell}).
\label{eq:semantic_loss}
\end{equation}
The target serialization is fixed, for example,
\begin{equation}
\texttt{<INTENT=i;ATTRIBUTE=a;CONTEXT=c>},
\label{eq:serialization_example}
\end{equation}
and decoding accepts only tokens in the corresponding field vocabularies.
An output is valid if and only if it contains exactly one value for every
field and every value belongs to the predefined vocabulary.
Invalid outputs are regenerated once under greedy decoding; a second invalid
output is mapped to an explicit \texttt{UNK} code and counted in the reported
invalid-output rate.

To prevent trajectory-label leakage, $G_{\phi}$ receives POI metadata only.
It does not receive user IDs, check-in sequences, target transitions, or
validation/test labels.
In the standard transductive setting, metadata of the catalog may be encoded
offline, but all interaction statistics used by the retrieval experts are
estimated exclusively from the training split.
Any cold-POI claim requires a separate inductive protocol in which the
semantic generator is not tuned on the held-out POIs.

\subsubsection{Retrieval Representation}

The implementation serializes the structured identifier together with POI
metadata and encodes the resulting text with the retrieval encoder:
\begin{equation}
\bm{e}_{p}^{\mathrm{gs}}
=
\operatorname{Enc}_{\mathrm{ret}}
\left(
\operatorname{Serialize}
\left[
\gsid(p),\bm{m}_p
\right]
\right).
\label{eq:geosemid_embedding}
\end{equation}
The serialization retains the canonical POI ID for exact identity and exposes
all GeoSemID fields to the same encoder used for trajectory retrieval. This
matches both local and API embedding backends and does not assume separately
trained identifier embedding tables or an additional projection layer.

\subsection{Reliability-Aware Hierarchical Backoff for Count Statistics}
\label{subsec:backoff}

Transition and temporal probabilities are estimated from finite interaction
counts. For either count-based expert $m\in\{\mathrm{tr},\mathrm{tp}\}$, let
$P_{u,\mathrm{ex}}^{(m)}(p)$ be the exact POI-level probability and
$P_{u,l}^{(m)}(p)$ be the POI-level probability induced by identifier level
$l$. The support of the exact statistic is $N_u^{(m)}$, and its reliability is
\begin{equation}
\lambda_{u}^{(m)}
=
\frac{N_{u}^{(m)}}{N_{u}^{(m)}+\alpha_m},
\label{eq:reliability}
\end{equation}
where $\alpha_m>0$ is fixed on the validation split.
The probabilities are interpolated before taking the logarithm:
\begin{align}
P_u^{(m)}(p)
={}&
\lambda_u^{(m)}P_{u,\mathrm{ex}}^{(m)}(p)
\nonumber\\
&+
(1-\lambda_u^{(m)})
\sum_{l\in\mathcal{L}_m}
\omega_l^{(m)}P_{u,l}^{(m)}(p),
\label{eq:unified_backoff_probability}\\
r_u^{(m)}(p)
={}&
\log\!\left(P_u^{(m)}(p)+\epsilon\right),
\label{eq:unified_backoff}
\end{align}
Here $\epsilon>0$ is a numerical safeguard,
$\omega_l^{(m)}\ge0$, and $\sum_l\omega_l^{(m)}=1$.
For an identifier value $z$, define its POI-bucket size as
\begin{equation}
\nu_l(z)
=
\left|\{p'\in\POIs:z_{p'}^{(l)}=z\}\right|.
\label{eq:identifier_bucket_size}
\end{equation}
Dividing an identifier-level probability by $\nu_l(z)$ below turns it into a
distribution over individual POIs rather than assigning the full code
probability to every POI in the bucket.
The smoothing constants, $\alpha_m$, and $\omega_l^{(m)}$ are selected on the
validation split using candidate Recall@$K_0$ and then fixed for test-time
inference. Spatial distance and semantic similarity are not count estimates;
their calibrated fusion is defined separately below.

For count probabilities, additive smoothing is uniquely defined as
\begin{equation}
P_{\delta}(b\mid a)
=
\frac{N(a,b)+\delta}{N(a)+\delta|\mathcal{Y}|},
\label{eq:additive_smoothing}
\end{equation}
where $\mathcal{Y}$ is the corresponding outcome vocabulary.
No probability term is silently omitted.

\subsection{Specialized Retrieval Experts}
\label{subsec:experts}

\subsubsection{Spatial Expert}

Let $\bm{g}_{u}^{\mathrm{last}}$ be the latest coordinate and
$\bm{g}_{u}^{\mathrm{ctr}}$ be the median center of the current trajectory.
The distance-based component is mapped to $[0,1]$ as
\begin{equation}
\begin{aligned}
s_{u,\mathrm{dist}}^{(\mathrm{sp})}(p)
={}&\frac{1}{2}\exp\!\left(-\widetilde d_{\mathrm{hav}}
(\bm{g}_{u}^{\mathrm{last}},\bm{g}_{p})\right)\\
&+\frac{1}{2}\exp\!\left(-\widetilde d_{\mathrm{hav}}
(\bm{g}_{u}^{\mathrm{ctr}},\bm{g}_{p})\right),
\end{aligned}
\label{eq:spatial_exact}
\end{equation}
where $\widetilde d_{\mathrm{hav}}$ is Haversine distance divided by the
training-set median distance.
The hierarchical terms are
\begin{align}
b_{u,\mathrm{cg}}^{(\mathrm{sp})}(p)
&=
\Ind[z_p^{\mathrm{cg}}\in\mathcal{R}_{u}^{\mathrm{cg}}],
&
b_{u,\mathrm{fg}}^{(\mathrm{sp})}(p)
&=
\Ind[z_p^{\mathrm{fg}}\in\mathcal{R}_{u}^{\mathrm{fg}}],
\label{eq:spatial_backoff}
\end{align}
where the active-region sets are estimated from the observable trajectory.
The spatial score uses the same $[0,1]$ scale for distance and region evidence:
\begin{equation}
r_u^{(\mathrm{sp})}(p)
=
\zeta_{\mathrm{sp}}s_{u,\mathrm{dist}}^{(\mathrm{sp})}(p)
+
(1-\zeta_{\mathrm{sp}})
\sum_{l\in\{\mathrm{cg},\mathrm{fg}\}}
\omega_l^{(\mathrm{sp})}b_{u,l}^{(\mathrm{sp})}(p),
\label{eq:spatial_fusion}
\end{equation}
where $\zeta_{\mathrm{sp}}\in[0,1]$ and $\omega_l^{(\mathrm{sp})}$ are selected
on the validation split by candidate Recall@$K_0$.

\subsubsection{Transition Expert}

Let $p_u^{\mathrm{last}}$ be the latest POI.
The exact transition probability is
\begin{equation}
P_{u,\mathrm{ex}}^{(\mathrm{tr})}(p)
=
P_{\delta_{\mathrm{tr}}}(p\mid p_u^{\mathrm{last}}).
\label{eq:transition_exact}
\end{equation}
For $l\in\{\mathrm{fg},\mathrm{cat},\mathrm{sem}\}$,
\begin{equation}
P_{u,l}^{(\mathrm{tr})}(p)
=
\frac{
P_{\delta_{\mathrm{tr}},l}
(z_p^{(l)}\mid z_{p_u^{\mathrm{last}}}^{(l)})
}{
\nu_l(z_p^{(l)})
}.
\label{eq:transition_backoff}
\end{equation}
The support $N_u^{(\mathrm{tr})}$ in \eqref{eq:reliability} is the number of
observed outgoing transitions from $p_u^{\mathrm{last}}$ in the training
split; it is shared by all candidate POIs for the same prediction instance.

\subsubsection{Temporal Expert}

Let $\tau_u$ be the joint time-of-day, day-of-week, and weekday/weekend bin.
The exact probability is
\begin{equation}
P_{u,\mathrm{ex}}^{(\mathrm{tp})}(p)
=
P_{\delta_{\mathrm{tp}}}(p\mid u,\tau_u).
\label{eq:temporal_exact}
\end{equation}
The backoff levels are
\begin{equation}
P_{u,l}^{(\mathrm{tp})}(p)
=
\frac{P_{\delta_{\mathrm{tp}},l}(z_p^{(l)}\mid\tau_u)}
{\nu_l(z_p^{(l)})},
\quad
l\in\{\mathrm{cat},\mathrm{sem}\}.
\label{eq:temporal_backoff}
\end{equation}
Its exact support is the number of user-$u$ check-ins observed in bin
$\tau_u$; low support automatically shifts mass to category and semantic
statistics through \eqref{eq:reliability}.

\subsubsection{Semantic Expert}

The retrieval encoder produces an intent query from the serialized observable
trajectory:
\begin{equation}
\bm{q}_{u}^{\mathrm{sem}}
=
\operatorname{Enc}_{\mathrm{ret}}
\left(
\operatorname{Serialize}
\left[\mathcal{T}_{u}^{H},\mathcal{T}_{u}^{C}\right]
\right).
\label{eq:semantic_query}
\end{equation}
The semantic-similarity component is normalized to $[0,1]$:
\begin{equation}
s_{u,\mathrm{sim}}^{(\mathrm{se})}(p)
=
\frac{1+\operatorname{cos}\!\left(
\bm{q}_{u}^{\mathrm{sem}},\bm{e}_{p}^{\mathrm{gs}}
\right)}{2},
\label{eq:semantic_exact}
\end{equation}
Its identifier-level evidence is
\begin{equation}
\begin{aligned}
b_{u,l}^{(\mathrm{se})}(p)
&=
\frac{1}{|\mathcal{T}_{u}^{C}|}
\sum_{x_{u,i}\in\mathcal{T}_{u}^{C}}
\Ind[z_{p_{u,i}}^{(l)}=z_p^{(l)}],\\
&\hspace{1.2em}
 l\in\{\mathrm{cg},\mathrm{fg},\mathrm{cat},\mathrm{sem}\}.
\end{aligned}
\label{eq:semantic_backoff}
\end{equation}
The semantic score combines two quantities on the same $[0,1]$ scale:
\begin{equation}
r_u^{(\mathrm{se})}(p)
=
\zeta_{\mathrm{se}}s_{u,\mathrm{sim}}^{(\mathrm{se})}(p)
+
(1-\zeta_{\mathrm{se}})
\sum_{l\in\{\mathrm{cg},\mathrm{fg},\mathrm{cat},\mathrm{sem}\}}
\omega_l^{(\mathrm{se})}b_{u,l}^{(\mathrm{se})}(p),
\label{eq:semantic_fusion}
\end{equation}
where $\zeta_{\mathrm{se}}\in[0,1]$ and $\omega_l^{(\mathrm{se})}$ are selected
on the validation split by candidate Recall@$K_0$. This formulation makes the
contribution of every GeoSemID component observable and directly ablatable.

\subsection{Utility-Supervised Sample-Level Router}
\label{subsec:router}

\subsubsection{Observable Routing Context}

The router receives no target information.
Its input is
\begin{equation}
\bm{q}_{u}
=
[
\bm{h}_{u}^{H}
\Vert
\bm{h}_{u}^{C}
\Vert
\bm{e}_{\tau_u}
\Vert
\bm{s}_{u}^{\mathrm{mob}}
\Vert
\bm{s}_{u}^{\mathrm{gs}}
],
\label{eq:routing_context}
\end{equation}
where $\bm{s}_{u}^{\mathrm{mob}}$ includes displacement, radius of gyration,
and trajectory concentration, and $\bm{s}_{u}^{\mathrm{gs}}$ summarizes
GeoSemID frequencies and transitions in the observable trajectory.
The router predicts
\begin{equation}
\bm{\pi}_{u}
=
\operatorname{softmax}(f_{\theta}(\bm{q}_{u})/T_r).
\label{eq:router_output}
\end{equation}

\subsubsection{Training-Only Rank-Derived Utility}

During router training only, let $\operatorname{rank}_{u}^{(m)}(p_u^{*})$ be
the target rank in expert $m$'s top-$K_0$ list when the target is retrieved.
The utility is
\begin{equation}
a_{u}^{(m)}
=
\begin{cases}
1-\dfrac{\operatorname{rank}_{u}^{(m)}(p_u^{*})-1}{K_0},
&p_u^{*}\in\mathcal{C}_{u}^{(m)},\\[7pt]
0,&p_u^{*}\notin\mathcal{C}_{u}^{(m)},
\end{cases}
\label{eq:expert_utility}
\end{equation}
and the soft target is
\begin{equation}
q_{u}^{(m)}
=
\frac{
\Ind[a_u^{(m)}>0]\exp(a_{u}^{(m)}/\tau_q)
}{
\sum_{j\in\Experts}
\Ind[a_u^{(j)}>0]\exp(a_{u}^{(j)}/\tau_q)
}.
\label{eq:utility_target}
\end{equation}
We train the router only on examples that are retrieved by at least one expert:
\begin{equation}
\mathcal{D}_{\mathrm{router}}
=
\{u\in\Train:\max_{m\in\Experts}a_u^{(m)}>0\}.
\label{eq:router_training_set}
\end{equation}
This aligns supervision with the top-$K_0$ candidate-generation task: an expert
is rewarded only when it can actually contribute the target to the union pool.
The router loss is
\begin{equation}
\mathcal{L}_{\mathrm{router}}
=
-
\sum_{u\in\mathcal{D}_{\mathrm{router}}}
\sum_{m\in\Experts}
q_{u}^{(m)}\log\pi_{u}^{(m)}.
\label{eq:router_loss}
\end{equation}
All utilities are generated from training targets only.
Validation and test targets are never used to compute router weights or
routing features.

\subsection{Comparable Dynamic Fusion and Evidence Preservation}
\label{subsec:fusion}

Raw expert scores have incompatible scales, so they are not added directly.
Within each top-$K_0$ expert list, we convert rank to a percentile score:
\begin{equation}
\overline r_{u}^{(m)}(p)
=
\begin{cases}
1-\dfrac{\operatorname{rank}_{u}^{(m)}(p)-1}{K_0},
&
p\in\mathcal{C}_{u}^{(m)},\\[7pt]
0,
&
p\notin\mathcal{C}_{u}^{(m)}.
\end{cases}
\label{eq:rank_percentile}
\end{equation}
The union pool and dynamically fused score are
\begin{align}
\widetilde{\mathcal{C}}_u
&=
\bigcup_{m\in\Experts}\mathcal{C}_{u}^{(m)},
\label{eq:union_pool}\\
s_u(p)
&=
\sum_{m\in\Experts}
\pi_u^{(m)}\overline r_u^{(m)}(p).
\label{eq:dynamic_fusion}
\end{align}
The final candidate set is
\begin{equation}
\mathcal{C}_u
=
\operatorname{TopK}_{K}
\{s_u(p):p\in\widetilde{\mathcal{C}}_u\}.
\label{eq:final_candidate}
\end{equation}
Ties are broken first by the number of experts retrieving the candidate and
then by canonical POI ID.

For every retained candidate, the evidence vector is
\begin{equation}
\bm{\xi}_{u,p}
=
\mathop{\Vert}_{m\in\Experts}
[
\pi_u^{(m)},
\overline r_u^{(m)}(p),
\Ind[p\in\mathcal{C}_u^{(m)}]
].
\label{eq:evidence_vector}
\end{equation}
This vector is passed to the ranker; retrieval evidence is not discarded after
fusion.
Fixed reciprocal-rank fusion is not part of the proposed method and is used
only as a non-learned comparison in experiments.

\subsection{Reproducible Candidate-Constrained LLM Ranking}
\label{subsec:ranker}

\subsubsection{Structured Prompt}

Each training and inference prompt uses the following fixed fields:
\begin{enumerate}[leftmargin=*,nosep]
\item \texttt{HISTORY\_PROFILE}: a serialized summary derived only from
$\mathcal{T}_{u}^{H}$;
\item \texttt{CURRENT\_TRAJECTORY}: the last $L$ check-ins with timestamps;
\item \texttt{PREDICTION\_TIME}: the target time context available at inference;
\item \texttt{CANDIDATES}: one record per candidate containing canonical POI
ID, GeoSemID fields, metadata, router weight, expert memberships, and
rank-percentile scores;
\item \texttt{OUTPUT}: exactly one canonical candidate POI ID.
\end{enumerate}
Candidate records are randomly permuted at every training epoch to reduce
position memorization.
At inference, candidates are ordered by canonical POI ID, not by fused score.
Inputs longer than the fixed maximum length $L_{\max}$ retain all candidate
IDs and evidence while truncating the oldest historical records first.

\subsubsection{Explicit Candidate Likelihood}

Let $\operatorname{tok}(p)$ be the tokenizer sequence of the canonical POI ID.
The model does not freely generate an arbitrary catalog ID.
Instead, all candidate strings are scored by batched teacher forcing:
\begin{equation}
\psi_u(p)
=
\frac{1}{|\operatorname{tok}(p)|}
\sum_{\ell=1}^{|\operatorname{tok}(p)|}
\log
P_{\psi}
\left(
\operatorname{tok}_{\ell}(p)
\mid
\bm{x}_u,
\operatorname{tok}_{<\ell}(p)
\right).
\label{eq:candidate_likelihood}
\end{equation}
The candidate probability is
\begin{equation}
P_{\psi}(p\mid\bm{x}_u,\mathcal{C}_u)
=
\frac{\exp(\psi_u(p)/T_p)}
{\sum_{p'\in\mathcal{C}_u}\exp(\psi_u(p')/T_p)}.
\label{eq:candidate_probability}
\end{equation}
This protocol has decoder cost
$O(K\overline{\ell}_{\mathrm{id}})$ token evaluations per sample and can be
batched across candidates.
It also makes candidate order, ID tokenization, and output constraints
explicitly reproducible.

\subsection{Training and Inference Consistency}
\label{subsec:training}

Training proceeds in three stages.

\textit{Stage 1: semantic-code generator tuning.}
The Qwen-8B generator is optimized with
$\mathcal{L}_{\mathrm{sem}}$, validated under the fixed grammar, and used once
to cache $z_p^{\mathrm{sem}}$ for every eligible POI.

\textit{Stage 2: expert construction and router training.}
All transition and temporal counts, spatial statistics, and semantic indexes
are built from the training interaction split.
Top-$K_0$ training retrieval lists provide utility targets for
$\mathcal{L}_{\mathrm{router}}$.

\textit{Stage 3: candidate-constrained ranker tuning.}
The trained retrievers and router generate the natural candidate set
$\mathcal{C}_u$.
No true target is inserted.
The ranker training subset is
\begin{equation}
\HitTrain
=
\{u\in\Train:p_u^{*}\in\mathcal{C}_u\},
\label{eq:hit_training_set}
\end{equation}
and the ranking loss is
\begin{equation}
\mathcal{L}_{\mathrm{rank}}
=
-
\sum_{u\in\HitTrain}
\log
P_{\psi}(p_u^{*}\mid\bm{x}_u,\mathcal{C}_u).
\label{eq:rank_loss}
\end{equation}
Missed samples are not made artificially easier by target injection; instead,
they remain retrieval errors and are addressed through the expert and router
stages.

At inference, all three online steps use only observable information:
expert retrieval, dynamic fusion, and candidate likelihood scoring.
The semantic generator is not invoked online.

\subsection{Evaluation Contract}
\label{subsec:evaluation_contract}

To prevent candidate recall and ranking quality from being conflated, the
experiments must report:
\begin{enumerate}[leftmargin=*,nosep]
\item candidate Recall@$K$ on the unmodified candidate set;
\item ranker-training coverage
$|\mathcal{D}_{\mathrm{hit}}|/|\mathcal{D}_{\mathrm{tr}}|$ and the corresponding
number of ranker-training samples;
\item conditional ranking HR/NDCG/MRR on samples where
$p_u^{*}\in\mathcal{C}_u$;
\item end-to-end HR/NDCG/MRR, counting every retrieval miss as incorrect;
\item an explicitly labeled oracle-candidate upper bound, used only for
diagnosis and never as the main result;
\item long-tail POI, short-trajectory, and sparse-user slices;
\item offline GeoSemID generation cost, online expert retrieval latency,
router latency, and candidate-ranking latency.
\end{enumerate}

The minimum mechanism ablations are:
GeoSemID removed; geographic/category codes only; semantic code only;
semantic embedding without discrete codes; uniform fusion; globally fixed
weights; fixed RRF; removal of each expert; removal of router utility
supervision; and removal of expert evidence from the ranker.

\bibliographystyle{IEEEtran}
% \bibliography{references}

\end{document}
